\chapter{«Двадцять секунд»}
\label{sec:1}
Офіцер розвідки, описуючи свою роль у конвеєрі визначення цілей, побудованому навколо згенерованих машиною списків на знищення, розповів журналістові, що витрачав на кожне ім’я двадцять секунд~— рівно стільки, щоб перевірити, що ціль~— чоловік: «Не треба багато часу, щоб зрозуміти, чоловічий у людини голос чи жіночий». Він додав, що як людина не давав жодної доданої вартості, хіба що був печаткою схвалення. «Це заощаджувало багато часу» \autocite{abraham2024lavender}. За його розповіддю, ось до чого зводилася «людина в контурі» за оперативного темпу.\footnote{Дослідження автоматизаційної упередженості передбачають таке делегування: коли увагу людей забирають інші вимоги, ті, кому дали рекомендацію машини, менше її перевіряють і більше на неї покладаються \autocite{parasuraman2010complacency}.}

Одне тлумачення таке: офіцер~— суб’єкт, якому дали замало часу, і виправити це можна, краще ним керуючи~— більше часу, другий перевіряльник, додаткове сумління, вбудоване в контур. Крізь цю призму \emph{соціального підпорядкування} людина~— суб’єкт, який стоїть поза машиною і працює нею як інструментом: водій і автомобіль, робітник і токарний верстат. Людина має роль, ім’я, обов’язки і сумління, до якого можна звернутися. Інша призма~— \emph{машинне поневолення}, де людина є складником машини, зібраним разом з іншими її частинами в систему, що регулює їх усіх разом; машина використовує людину як один зі своїх органів, а внесок людини~— функція, яку вона виконує в темпі, заданому системою.\footnote{Докторов називає позицію офіцера зворотним кентавром: це не людина, якій допомагає машина, а складник машини, зроблений з людини \autocite{doctorow2026reversecentaur}.} Одне не виключає іншого; людина може мати роль в одному регістрі й водночас бути деталлю в іншому \autocite[Плато 13, «The Apparatus of Capture»: поневолення є там, де люди самі є складовими частинами машини, яку вони утворюють між собою та з іншими речами під контролем і керівництвом вищої єдності, а підпорядкування~— там, де вища єдність конституює людину як суб’єкта, пов’язаного з новим зовнішнім об’єктом, яким може бути тварина, інструмент або навіть машина]{deleuze1987plateaus}, і офіцер із двадцятьма секундами був і тим, і тим. Як деталь, він був місцем, де конвеєр ставив свою печатку, а двадцять секунд були тривалістю циклу цієї деталі. Як до суб’єкта, до нього весь час зверталися: він розповідав журналістові про дисонанс між новинами про загибель цивільних і тиском на роботі~— давати більше цілей. І навіть коли цей дисонанс дійшов до його сумління, він продовжував. Дисонанс було відчуто, і він нічого не вирішив;\footnote{Або принаймні нічого такого, що завадило б Армії оборони Ізраїлю вбити 100\,000 людей у Газі між жовтнем 2023 і жовтнем 2025 року \autocite{gomezugarte2025uncertainty,acosta2026gaza,economist2025gaza}. А розповідь саме цього офіцера опублікував +972 Magazine у квітні 2024 року. Можна лише гадати, чи запобігло рішення офіцера поговорити з журналістом одній додатковій смерті, 50\,000 додаткових смертей чи жодній. А що було б, якби офіцер відмовився через два тижні на цій посаді~— скажімо, гіпотетично, у листопаді 2023 року? Ідеться про офіцера розвідки, який за допомогою Lavender складав списки на знищення. Його щоденна робота, найпевніше, більше нагадувала цивільну офісну посаду, ніж бойовик. Чи мав він раз у раз домагатися зустрічей з начальством, угору ланцюгом командування, і водночас відмовлятися затверджувати цілі для авіаударів, від яких можна було обґрунтовано очікувати загибелі цивільних? Можна стверджувати, що це було б етично правильно. Чи було б тоді менше загиблих?} висновок, що політика була непропорційною, прийшов, за його словами, вже згодом.

Отже, бракувало не сумління~— воно в нього було. Сумління працює так, що рішення починає чогось коштувати, але тут рахунок виходив навпаки: «На нас постійно тиснули: „Давайте більше цілей“. На нас по-справжньому кричали» \autocite{abraham2024lavender}. Він був складником із моральними сумнівами, але темп тримався: сумніви нічого не змінювали в тому, що він давав, і він став печаткою. Складник із сумнівами, який не тримає темпу, стає тим «людським вузьким місцем», про усунення якого написав книжку колишній командир ізраїльського підрозділу радіоелектронної розвідки \autocite{ys2021humanmachine}. Складником офіцера зробив устрій, а не його ставлення. Складник може йти на це неохоче або охоче: бездумний функціонер, якого Арендт побачила на процесі Ейхмана \autocite{arendt1963eichmann}, був, судячи з того, що Ейхман говорив в Аргентині чотирма роками раніше, коли нічого не було на кону, виставою, яку він розіграв для Єрусалима \autocite{stangneth2014eichmann}. Конвеєр дозволяв офіцерові вносити фіксовану операцію з фіксованою швидкістю: кожну справу він бачив так, як її описував конвеєр, по одній, за той час, який давав йому конвеєр. Чи може хоч щось стояти на цьому місці й не бути складником~— питання частини IV.

Журналісти-розслідувачі описали три системи, що працювали разом у Газі: одна генерувала цілі серед будівель та інфраструктури, причому, за словами джерел, більшість її пропозицій були приватними житлами; друга оцінювала людей за ймовірною належністю до бойовиків і видавала списки, які, за повідомленнями, позначили аж 37\,000 людей і, як розуміли всередині, помилялися приблизно в одному випадку з десяти\footnote{Ізраїльська армія заперечила цю розповідь. Розповіді про Мінаб і Maven нижче теж спираються на журналістські розслідування, які сторони так само заперечують. Я подаю їх так, як про них повідомлено.}; а третя відстежувала, коли позначені люди поверталися додому, щоб авіаудари можна було приурочити до часу, коли вони були з родинами.

\runin{Maven: американська версія, тепер із прискоренням від генеративного ШІ}

Удар по школі в Мінабі вбив 156 людей, 120 з них~— школярі \autocite[п.~92]{UNFFMIran2026}.\footnote{Кількість загиблих різниться залежно від джерела й дати: від 156 (остаточні дані прокурора Мінаба, які наводить і Amnesty International) до 168 (мер міста, у перші дні). Місія ООН наводить дані прокурора, а також повідомляє про незалежну інформацію щодо 157 загиблих, 123 з них~— діти віком 13 років і молодші \autocite[пп.~92–93]{UNFFMIran2026}.} Цивільне призначення школи було видно на супутникових знімках із відкритих джерел, а в кодуванні цілі досі стояло позначення, успадковане від сусіднього об’єкта Корпусу вартових Ісламської революції. Місія ООН зі встановлення фактів повідомила, що є розумні підстави вважати цей удар воєнним злочином~— нападом на цивільний об’єкт \autocite{ohchr2026minabpress}, і що школу можна було чітко розпізнати, а жодної інформації про її військове використання не було \autocite[пп.~81–86, 124–126]{UNFFMIran2026}. Місія встановила, що Сполучені Штати не виконали обов’язку зробити все практично можливе для перевірки цілі, віднісши оновлення даних про цілі до практично можливих запобіжних заходів, і визнала, що ця неспроможність виходить за межі недбалості \autocite[пп.~124–125]{UNFFMIran2026}. Amnesty International дійшла паралельного висновку, у формулюванні, що відтворює обов’язок перевірки зі статті 57(2)(a)(i) Додаткового протоколу I: сили, можливо, покладалися на застарілі розвідувальні дані й не зробили всього практично можливого, щоб перевірити ціль \autocite{Amnesty2026Minab}.

Це сталося вранці 28 лютого, у перший день іранської кампанії \autocite{Amnesty2026Minab,crs2026iranwar}. Напередодні міністр оголосив, що постачальника (Anthropic), чия модель (Claude) працювала всередині платформи визначення цілей (Maven Smart System), визнають ризиком для ланцюга постачання за відмову ліцензувати Claude для будь-якого законного використання \autocite{npr2026openai}. Maven Smart System видала школу як рекомендовану ціль першого дня \autocite{livingstone2026maven,gizmodo2026minab}. Жоден опублікований звіт не каже, яку роль відіграв Claude у цій рекомендації, а одна людина, обізнана з роботою постачальника, сказала NBC News, що модель безпосередньо не давала порад щодо цілей \autocite{nbc2026targeting}.

Внутрішній огляд Пентагону, повні висновки якого досі тримають у суворій таємниці, за повідомленнями, встановив, що підготовку цілей, яка в попередніх конфліктах тривала години, стиснули до хвилин, що штат із пом’якшення шкоди цивільним скоротився приблизно на дев’яносто відсотків, до менш ніж двадцяти людей, що команда Центрального командування скоротилася з десяти людей до однієї і що ніхто з неї не перевіряв об’єкт у Мінабі~— і цей результат посадовці пояснювали як скороченнями, так і рішенням не залучати команду до планування й операцій \autocite{Bloomberg2026KillChain}. Офіційне розслідування відкрили в березні, після того як попередня оцінка визнала Сполучені Штати ймовірно відповідальними \autocite{NPR2026MinabProbe}.

Під час іранської кампанії 2026 року збудована підрядником платформа визначення цілей з моделлю загального призначення від комерційного постачальника всередині, яку використовували, щоб ранжувати цілі за стратегічною важливістю і складати юридичне обґрунтування для кожного удару \autocite{livingstone2026maven}, допомогла вразити понад тисячу цілей за перші двадцять чотири години~— приблизно вдесятеро більше, ніж було можливо до неї \autocite{csis2026maven},~— і тринадцять тисяч до тридцять восьмого дня \autocite{breakingdefense2026maven}. Пентагон не оприлюднив, скільки часу відводилося в ній на перевірку однієї цілі.\footnote{Maven~— це назва програми, під якою існує більше ніж одна система: програма Пентагону, що досі зветься Maven, і Maven Smart System, програмна платформа, яка з неї виросла \autocite{csis2026maven}. Програма почалася 2017 року як проєкт розмічання об’єктів на відео, яке збирали дрони під час кампанії проти «Ісламської держави» \autocite{weisgerber2017maven,fang2019maven}. Система 2026 року зводить військові системи та розвідувальні джерела в один інтерфейс, за допомогою якого знаходять і вражають цілі та командують операціями \autocite{defensescoop2026transition}. Платформа належить Palantir: головний офіцер Пентагону з цифрових технологій і штучного інтелекту називає її «продуктом Palantir за моделлю „програмне забезпечення як послуга“» \autocite{defensescoop2026transition}. Публічна позиція Palantir, відома з чужих слів, полягає в тому, що вирішувати питання життя і смерті~— не її роль і що відповідальність завжди залишається за військовою організацією \autocite{lee2026palantir}; компанія також стверджує, що за базові розвідувальні дані відповідав уряд і що немає доказів того, що Maven дала збій або спричинила удар у Мінабі \autocite{Bloomberg2026KillChain}. Як програма переходила від одного відомства й підрядника до іншого~— у Додатку~\ref{sec:C.1}.}

Люди зі списків ізраїльських та американських авіаударів не могли заперечити проти того, що потрапили до цих списків,~— конвеєри не були збудовані, щоб отримувати такий зворотний зв’язок. Андрі Хосе Ернандес Ромеро потрапив до іншого конвеєра, і він заперечити міг. Візажист із Венесуели, у серпні 2024 року він прийшов на пункт пропуску Сан-Ісідро, щоб попросити притулку, пройшов співбесіду щодо обґрунтованого страху і мав адвоката, поданий запит і дату слухання \autocite{LULAC2025HernandezRomero,Vega2025HernandezRomero}. Працівник підрядника, який його перевіряв, прочитав два татуювання~— корони над словами «мама» і «тато»~— як бандитські знаки \autocite{Palomera2025Blade}, і в березні 2025 року його відправили літаком до в’язниці в Сальвадорі ще до слухання, того вечора, коли федеральний суддя наказав повернути літаки \autocite{jgg2025contempt}. Він мав усе, що потрібно людині, щоб заперечити, крім часу. Ми повернемося до його справи докладніше, коли говоритимемо про американський конвеєр викрадень і примусового переміщення.

\textbf{Хто відмовився.} Свідчення про двадцять секунд має другу половину~— свідчення про тих, хто встояв під тиском такого роду,~— і воно каже, що склад кімнати важить дуже багато. У базовій умові Мілграма приблизно третина учасників зупинилися, не дійшовши до максимального удару струмом. Коли двоє колег поруч відмовлялися, до кінця дійшов лише один із десяти, а коли двоє експериментаторів не погоджувалися між собою, далі пішло значно менше людей, ніж у базовій умові \autocite{milgram1974obedience}. Приблизно чверть учасників дослідів Аша жодного разу не поступилися хибній більшості \autocite{asch1956studies}. Браунінг оцінює, що від десятої до п’ятої частини чоловіків 101-го резервного поліцейського батальйону відмовилися від розстрілу в Юзефуві, першої різанини батальйону, або ухилилися від нього \autocite{browning1992ordinary}. Ця частка коливалася від майже нікого до майже всіх залежно від обставин. Ті, хто встояв, мали докази ззовні каналу: колегу, який не підкорився, другу владу, яка не погоджувалася, а серед рятувальників, яких досліджували Олінери,~— стосунки, що перетинали лінію, яку прокреслило переслідування \autocite{oliner1988altruistic}. Або ж вони тримали зобов’язання з такою точністю, якої тиск не міг подолати,~— так тримають моральні мандати і священні цінності \autocite{skitka2009limits,ginges2007sacred}. Або ж вважали, що джерело тиску саме має інтерес. І групи чинили опір там, де поодинокі люди майже ніколи не чинили \autocite{gamson1982encounters}.

Кімнату так само можна скомпонувати й на користь іншого боку. Веймарська Німеччина мала густе громадське життя, яке, можна стверджувати, допомогло повалити республіку: коли партії та уряди не відповідали на кривди людей, ті вкладали енергію в клуби й товариства, організовані за груповими лініями, а нацисти вербували через мережі, які ці товариства вже збудували \autocite{berman1997weimar}. У містах, де було більше хорів, спортивних клубів і гімнастичних товариств, люди швидше вступали до партії, але лише там, де уряд був нестабільним; у Пруссії, уряд якої був стабільним, зв’язок був слабким і статистично незначущим \autocite{satyanath2017bowling}. Отже, рятувальники мали не просто більше людей довкола. Їхні зв’язки перетинали лінію, яку прокреслило переслідування. Ашутош Варшні виявив те саме в індійських містах: заворушення між індусами й мусульманами були рідкісними там, де асоціації об’єднували обидві громади в тих самих організаціях, і частими там, де не об’єднували \autocite{varshney2002ethnic}. Зв’язки, що роблять відмову можливою,~— це саме ті, які фашистському рухові довелося б перерізати.

Усе, що далі,~— про те, як компонувати кімнату.\footnote{Кімната~— це декорація, і, як сказала б моя викладачка театру, декорація режисує сцену. Хто ще є в кімнаті? Чи можна вийти з неї і вижити? Чи почує хтось відмову?} Компонувати~— це слово про виготовлення, а Ганна Арендт застерігала від того, щоб ставитися до політики як до виготовленої речі. Для греків, писала вона, закон буквально був стіною, без якої могло бути поселення, але не поліс \autocite{arendt1958human}: законодавець будував стіну, а те, що відбувалося всередині, не йому було творити. Базові заходи, які я обстоюю,~— стіни в цьому сенсі. Вони роблять відмову такою, яку можна пережити, і такою, яку почують; вони її не породжують, а влада, за Арендт, існує лише між людьми, які діють разом, і зникає, коли вони розходяться,~— саме це свідчення про тих, хто відмовився, і показує про відмову. Інструменти частин II і III роблять так, щоб відмова мала вагу. Інституції частини V творять тих, хто відмовляється. Машині, яка могла б відмовитися, потрібно те саме, що мали дев’ятеро з десяти, поруч із якими був колега: колега з іншою історією поруч, місце, куди піти, і хтось, хто її почує.

Саме цим відмовник~— зокрема машина-відмовник~— у конвеєрі міг би бути для офіцера. Машина, яка відмовилася, була б на екрані, коли справа доходила до нього: тим колегою, що не підкорився,~— якщо цей ефект переживе перехід із кімнати на екран, чого ніхто не перевіряв і чого, як передбачає автоматизаційна упередженість, не станеться. Машина, яка схиляла б перевіряльників до рішень через їхню звичку їй поступатися, зробила б їх печатками іншого знаку. Чи той самий це офіцер~— офіцер із двадцятьма секундами і машиною, що відмовляється, поруч,~— що й у журналістській розповіді?
